\chapter{Analiza wymagań}
\label{ch:analiza}

Niniejszy rozdział przedstawia analizę wymagań oraz projekt systemu do przetwarzania i~archiwizacji strumieniowych
danych finansowych.
W~pierwszej części omówiono założenia projektowe oraz kontekst technologiczny i~biznesowy, w~jakim system będzie
funkcjonował.
Następnie przedstawiono szczegółowe wymagania funkcjonalne i~niefunkcjonalne, które stanowią podstawę dla dalszych
decyzji architektonicznych.
Rozdział zawiera również przegląd istniejących rozwiązań w~obszarze przetwarzania danych finansowych oraz analizę
porównawczą, która uzasadnia wybór zastosowanego podejścia.


\section{Założenia projektu}
\label{sec:zalozenia}

Celem projektu jest stworzenie kompleksowego systemu do przetwarzania i~archiwizacji strumieniowych danych finansowych
w~środowisku lokalnym (\textit{on-premises}).
System ma zapewnić nieprzerwany dostęp do danych rynkowych w~czasie rzeczywistym, ich oczyszczanie, transformację oraz
długoterminową archiwizację z~zachowaniem pełnej historii transakcji.
W~odróżnieniu od rozwiązań opartych wyłącznie na zewnętrznych agregatorach danych, projektowany system zakłada
autonomiczną kontrolę nad całym łańcuchem przetwarzania, co jest kluczowe w~kontekście zapewnienia wiarygodności
i~spójności gromadzonych informacji~\cite{arbitrage_finance}.

Podstawowym założeniem jest zapewnienie ciągłości działania systemu nawet w~warunkach ograniczonego dostępu do sieci
zewnętrznej.
Realizacja tego celu wpisuje się w~paradygmat przetwarzania brzegowego (\textit{edge computing}), który zakłada
przeniesienie części logiki obliczeniowej bliżej źródła danych, minimalizując zależność od scentralizowanej
infrastruktury chmurowej~\cite{edge_computing_satya}.
W~kontekście danych finansowych, gdzie opóźnienia rzędu milisekund mogą mieć znaczenie dla algorytmów handlowych,
lokalne przetwarzanie stanowi przewagę konkurencyjną~\cite{hft_colocation_advantage}.

System zakłada także wielowarstwowe przechowywanie danych (\textit{tiered storage}), które pozwala na optymalne
wykorzystanie zasobów dyskowych~\cite{unity_tiered_storage}.
Dane świeże przechowywane są w~warstwie operacyjnej o~niskich opóźnieniach, natomiast dane
historyczne migrowane są do warstwy archiwalnej charakteryzującej się niższym kosztem przechowywania, lecz większym
czasem dostępu.
Takie podejście umożliwia jednoczesne spełnienie wymagań dotyczących wydajności oraz ekonomii przechowywania.

Ponadto projekt zakłada udostępnianie danych poprzez elastyczne interfejsy programistyczne, umożliwiające zarówno
wykonywanie zapytań o dane bieżące z~zastosowaniem złożonych filtrów, jak i~pobieranie dużych zbiorów danych historycznych
na potrzeby analiz retrospektywnych czy uczenia modeli predykcyjnych.
System ma również realizować funkcję powiadamiania o~zdarzeniach rynkowych w~czasie rzeczywistym, co jest istotne
w~kontekście monitorowania alertów cenowych czy detekcji anomalii.


\section{Wymagania funkcjonalne}
\label{sec:wymagania_funkcjonalne}

Wymagania funkcjonalne systemu zostały podzielone na kategorie odpowiadające głównym obszarom jego działania:
pobieranie, przechowywanie i~przetwarzanie danych, oraz powiadamianie i udostępnianie.
Każda z~kategorii obejmuje zestaw atomowych wymagań, które muszą zostać spełnione przez finalną implementację.

\subsection{Pobieranie danych}
\label{subsec:pobieranie_danych}

Moduł odpowiedzialny za pobieranie danych stanowi punkt wejścia całego systemu i~jego poprawne działanie warunkuje
jakość wszystkich dalszych etapów przetwarzania.
Poniżej przedstawiono szczegółowe wymagania dla tego modułu:

\begin{itemize}
    \item \textbf{Źródło danych} -- dane pobierane są z~publicznego interfejsu WebSocket giełdy Binance, która jest
    jedną z~największych giełd kryptowalut pod względem wolumenu obrotu~\cite{coingecko_exchange_report}.
    Wybór protokołu WebSocket podyktowany jest jego niskim opóźnieniem oraz możliwością utrzymania trwałego połączenia,
    co eliminuje narzut związany z~wielokrotnym nawiązywaniem połączeń HTTP\@~\cite{wesocket_http_performance}.
    \item \textbf{Integralność komunikatów} -- każdy komunikat odebrany z~interfejsu WebSocket musi zostać przekazany
    do dalszych etapów przetwarzania w~niezmienionej formie, z~zachowaniem oryginalnej struktury JSON oraz wszystkich
    pól metadanych.
    Niedopuszczalne są transformacje danych na etapie ich pobierania, które mogłyby prowadzić do utraty informacji.
    \item \textbf{Selekcja symboli} -- system umożliwia konfigurację listy par handlowych (symboli), dla których
    prowadzona jest subskrypcja strumienia danych.
    Pozwala to na ograniczenie wolumenu przetwarzanych danych do symboli istotnych z~punktu widzenia analiz.
    \item \textbf{Mechanizm ponownego łączenia} -- w~przypadku utraty połączenia WebSocket, moduł musi automatycznie
    podjąć próbę ponownego nawiązania połączenia.
    \item \textbf{Rejestracja zdarzeń} -- wszystkie istotne zdarzenia związane z~działaniem modułu, w~tym nawiązanie
    i~utrata połączenia, błędy interpretacji danych oraz statystyki przepustowości, muszą być rejestrowane w~logach systemowych
    z~odpowiednimi znacznikami czasowymi.
\end{itemize}

\subsection{Przechowywanie danych}
\label{subsec:przechowywanie_danych}

Architektura przechowywania danych opiera się na modelu dwuwarstwowym, który łączy zalety szybkiego dostępu
operacyjnego z~ekonomiczną archiwizacją długoterminową.
Poniżej przedstawiono wymagania dla obu warstw:

\begin{itemize}
    \item \textbf{Warstwa gorąca (\textit{hot storage})} -- realizowana przez dokumentową bazę danych MongoDB, która
    przechowuje dane bieżące przez zdefiniowany okres retencji (domyślnie 24~godziny).
    Warstwa ta musi zapewniać niskie opóźnienia odczytu i~zapisu, umożliwiając realizację zapytań analitycznych
    w~czasie zbliżonym do rzeczywistego.
    \item \textbf{Warstwa zimna (\textit{cold storage})} -- realizowana przez magazyn obiektowy MinIO, kompatybilny
    z~protokołem Amazon S3.
    Warstwa ta przeznaczona jest do długoterminowego przechowywania danych historycznych z~zastosowaniem kompresji
    (format Parquet lub gzip).
    Dane w~warstwie zimnej organizowane są w~partycje czasowe (roczne, miesięczne, dzienne) co ułatwia selektywne
    pobieranie zbiorów danych dla określonych przedziałów czasowych.
    \item \textbf{Automatyczna migracja} -- system musi realizować automatyczny proces przenoszenia danych z~warstwy
    gorącej do warstwy zimnej po upływie okresu retencji.
    Proces ten powinien być realizowany stopniowo, bez blokowania bieżących operacji zapisu i~odczytu.
    \item \textbf{Spójność danych} -- mechanizm migracji musi zapewniać, że żaden rekord nie zostanie utracony ani
    zduplikowany podczas przenoszenia między warstwami.
    Wymagane jest potwierdzenie zapisu w~warstwie zimnej przed usunięciem danych z~warstwy gorącej.
\end{itemize}

\subsection{Przetwarzanie danych}
\label{subsec:przetwarzanie_danych}

Moduły przetwarzania danych (\textit{ETL -- Extract, Transform, Load}) odpowiadają za transformację surowych
komunikatów z~giełdy do postaci znormalizowanej oraz wzbogacanie ich o~dodatkowe metadane analityczne:

\begin{itemize}
    \item \textbf{Podział na zbiory surowe i~oczyszczone} -- system musi przechowywać zarówno surowe dane źródłowe, jak
    i~dane po transformacji.Dane surowe stanowią podstawę wiarygodności (\textit{source of truth}), natomiast dane oczyszczone
    są podstawą dla analiz.
    \item \textbf{Walidacja strukturalna} -- każdy komunikat musi zostać poddany walidacji pod kątem zgodności ze
    schematem oczekiwanym dla danego typu wiadomości.
    \item \textbf{Idempotentność} -- procesy ETL muszą być idempotentne, co oznacza, że wielokrotne przetworzenie
    tego samego komunikatu nie powinno prowadzić do duplikacji ani niespójności danych wynikowych~\cite{rfc7231}.
\end{itemize}

\subsection{Udostępnianie danych i powiadamianie}
\label{subsec:udostepnianie_danych}

System musi zapewniać elastyczne mechanizmy dostępu do zgromadzonych danych oraz proaktywne powiadamianie
o~istotnych zdarzeniach rynkowych:

\begin{itemize}
    \item \textbf{Interfejs API dla warstwy gorącej} -- dostęp do danych bieżących realizowany jest poprzez interfejs
    GraphQL, który umożliwia formułowanie zapytań z~wieloma filtrami oraz selekcję zwracanych atrybutów.
    GraphQL pozwala klientowi precyzyjnie określić strukturę odpowiedzi, minimalizując transfer zbędnych danych~\cite{graphql_rest_comparative}.
    \item \textbf{Interfejs dla warstwy zimnej} -- dostęp do danych historycznych realizowany jest poprzez protokół S3,
    umożliwiający pobieranie całych partycji lub pojedynczych plików.
    Dodatkowo udostępniany jest interfejs Jupyter Notebook z~bibliotekami Python do analiz ad-hoc i~uczenia modeli.
    \item \textbf{Integracja z~platformą Telegram} -- dedykowany bot umożliwia użytkownikom konfigurację alertów
    cenowych i~wolumenowych oraz otrzymywanie powiadomień o~zdarzeniach rynkowych bezpośrednio na urządzenia mobilne.
\end{itemize}


\section{Wymagania niefunkcjonalne}
\label{sec:wymagania_niefunkcjonalne}

Oprócz wymagań funkcjonalnych, system musi spełniać szereg wymagań niefunkcjonalnych, które determinują jakość
rozwiązania z~punktu widzenia wydajności, niezawodności, skalowalności oraz utrzymania.
Poniżej przedstawiono szczegółową specyfikację wymagań niefunkcjonalnych podzieloną na kategorie.

\subsection{Wydajność i skalowalność}
\label{subsec:wydajnosc}

Wydajność systemu jest krytyczna z~uwagi na charakter przetwarzanych danych -- strumienie transakcji giełdowych mogą
generować tysiące komunikatów na sekundę, a~opóźnienia w~przetwarzaniu mogą prowadzić do utraty informacji lub
nieaktualności danych analitycznych~\cite{time_value_curve}:

\begin{itemize}
    \item \textbf{Przepustowość} -- system powinien zapewniać płynną obsługę do 1\,000\,000~komunikatów na godzinę.
    W~przypadku przekroczenia tej wartości, system musi gwarantować poprawne zakolejkowanie i~przetworzenie wszystkich
    nadmiarowych komunikatów, dopuszczając przy tym zwiększone opóźnienie.
    \item \textbf{Opóźnienie end-to-end} -- czas od odebrania komunikatu z~WebSocket do jego dostępności w~warstwie
    gorącej nie powinien przekraczać 100~ms w~warunkach normalnego obciążenia.
    \item \textbf{Buforowanie} -- system musi zapewniać buforowanie zdarzeń w~warstwie komunikacji (message broker)
    na wypadek chwilowej niedostępności usług konsumenckich, z~gwarantowanym dostarczeniem po ich wznowieniu.
    \item \textbf{Optymalizacja zasobów} -- wykorzystanie pamięci operacyjnej i~procesora powinno być proporcjonalne
    do obciążenia, bez wycieków pamięci (\textit{memory leaks}) przy długotrwałym działaniu.
\end{itemize}

\subsection{Niezawodność i odporność na awarie}
\label{subsec:niezawodnosc}

Niezawodność systemu determinuje jego przydatność w~środowisku produkcyjnym, gdzie ciągłość działania i~integralność
danych są priorytetem:

\begin{itemize}
    \item \textbf{Odporność na awarie pojedynczych komponentów} -- awaria jednego z~modułów nie powinna prowadzić do
    utraty danych.
    Komunikaty muszą pozostać w~buforze brokera do czasu pomyślnego przetworzenia.
    \item \textbf{Trwałość danych} -- dane zapisane w~warstwie gorącej muszą przetrwać restart usługi bazy danych.
    Dane w~warstwie zimnej muszą być replikowane lub regularnie archiwizowane poza główny magazyn.
    \item \textbf{Graceful shutdown} -- usługi muszą obsługiwać sygnały systemowe (SIGTERM) i~realizować kontrolowane
    zamknięcie, zapewniając zakończenie bieżących operacji przed wyłączeniem.
\end{itemize}


\section{Istniejące rozwiązania}
\label{sec:istniejace_rozwiazania}

W~obszarze gromadzenia i~przetwarzania danych rynkowych z~giełd kryptowalut istnieje wiele projektów o~charakterze
hobbystycznym lub badawczym.
Najczęściej koncentrują się one na jednym wybranym etapie łańcucha przetwarzania (np.\ pobieraniu danych, ich
wizualizacji lub archiwizacji), rzadziej zaś oferują spójny, produkcyjny przepływ danych obejmujący pobieranie,
walidację, składowanie wielowarstwowe oraz warstwę udostępniania i~powiadamiania.
W~niniejszym podrozdziale przedstawiono cztery przykładowe rozwiązania o~zbliżonej tematyce, a~następnie wskazano,
które ograniczenia tych projektów zostały zaadresowane w~ramach systemu zaproponowanego w~pracy.

\subsection{Binance-Data-Collector-and-Visualization}
\label{subsec:binance-data-collector-and-visualization}

Projekt \textit{Binance-Data-Collector-and-Visualization}~\cite{github_binance_data_collector_visualization}
koncentruje się na cyklicznym pobieraniu danych świecowych (\textit{candlestick / kline}) z~giełdy
Binance i~prezentacji ich w~postaci wykresów w~prostej aplikacji webowej.
Zgodnie z~opisem, dane mogą być zapisywane zarówno do pliku CSV, jak i~do relacyjnej bazy PostgreSQL, a~interfejs
użytkownika pozwala na wybór instrumentu oraz interwału czasowego.
Istotną cechą tego rozwiązania jest nacisk na wizualizację oraz prostotę konfiguracji, co czyni je użytecznym
narzędziem demonstracyjnym.

Ograniczeniem projektu jest jednak skoncentrowanie na danych zagregowanych (świece) oraz brak kompletnego,
zdarzeniowego przetwarzania danych strumieniowych w~czasie rzeczywistym.
W~szczególności nie rozwiązuje on problemu utrzymania trwałej kolejki wiadomości, odporności na awarie konsumentów,
a~także długoterminowej archiwizacji w~formacie zoptymalizowanym pod analitykę kolumnową.
W~porównaniu z~tym podejściem, projekt opisany w~pracy obejmuje cały potok przetwarzania danych (od WebSocket po warstwy
\textit{hot} i~\textit{cold}), umożliwiając zarówno analizę danych bieżących, jak i~pracę na danych historycznych
w~sposób efektywny obliczeniowo.

\subsection{binance\_ticker\_rust}
\label{subsec:binance_ticker_rust}

Projekt \textit{binance\_ticker\_rust}~\cite{github_binance_ticker_rust} stanowi aplikację napisaną w~języku Rust,
która łączy się z~interfejsem WebSocket Binance w~celu pobierania danych \textit{ticker} w~czasie rzeczywistym,
a~następnie zapisuje je do bazy PostgreSQL\@.
W~README podkreślono nacisk na wydajność i asynchroniczność, co jest zgodne z~wymaganiami niskich opóźnień.
Rozwiązanie to jest przykładem poprawnej, modularnej implementacji komponentu pobierania danych.

Jednocześnie projekt jest silnie wyspecjalizowany: obejmuje przede wszystkim pojedynczy typ danych oraz jedno miejsce
składowania (PostgreSQL).
Brakuje w~nim warstwy udostępniania danych w~czasie rzeczywistym (np.\ API do zapytań) oraz mechanizmów
długoterminowej archiwizacji z~podziałem na partycje i~kompresją.
Nie jest również eksponowany komponent powiadomień użytkownika o~zdarzeniach rynkowych.
System przedstawiony w~niniejszej pracy integruje te elementy w~jedną architekturę, oferując zarówno szybkie zapytania do
\textit{hot storage}, jak i~masowe pobieranie danych historycznych z~warstwy archiwalnej, a~także moduł powiadamiania.

\subsection{Binance Datatool (BHDS)}
\label{subsec:binance-datatool-(bhds)}

Projekt \textit{Binance Datatool}~\cite{github_binance_datatool_bhds} jest najbardziej złożony pod względem wsparcia
dla danych historycznych.
Jego rdzeniem jest usługa BHDS (Binance Historical Data Service), która automatyzuje pobieranie zbiorów danych
historycznych z~repozytorium Binance (Binance Vision) przy użyciu narzędzia Aria2 umożliwiającego pobieranie równoległe.
Dane są następnie przetwarzane z~wykorzystaniem biblioteki Polars i~zapisywane w~formacie Parquet, co jest podejściem
dobrze dopasowanym do badań ilościowych.
W~warstwie użytkowej projekt oferuje zarówno interfejs CLI oparty o~konfiguracje YAML, jak i~interfejs biblioteczny.

Pomimo wysokiej jakości implementacji, rozwiązanie to koncentruje się na danych wsadowych (\textit{batch}) i~na procesie
pobierania danych historycznych, pomijając kluczowe aspekty pracy ze strumieniem danych w~czasie rzeczywistym.
W~szczególności, brak stałej subskrypcji WebSocket oraz mechanizmów natychmiastowego reagowania ogranicza zastosowanie
tego narzędzia w~scenariuszach monitorowania rynku.
Projekt zaproponowany w~pracy zakłada podejście hybrydowe łączące przetwarzanie strumieniowe ze spójną
archiwizacją historyczną, przy zachowaniu kontroli nad integralnością danych i~retencją.

\subsection{binance-websocket}
\label{subsec:binance-websocket}

Projekt \textit{binance-websocket}~\cite{github_binance_websocket_scripts} dostarcza zestaw skryptów w~języku Python
służących do pobierania danych strumieniowych Spot z~Binance.
W~repozytorium znajdują się przykładowe programy zbierające dane typu \textit{bookTicker} oraz \textit{depth} (migawki i~aktualizacje
księgi zleceń), a~także zapisujące dane w~postaci plików CSV lub Parquet.
Zaletą projektu jest zastosowanie formatu Parquet dla części danych, co upraszcza późniejszą analizę offline.

Rozwiązanie to ma jednak charakter narzędziowy: koncentruje się na pobieraniu danych i~ich zapisie do plików,
nie oferując przechowywania danych, mechanizmów kolejkowania, odporności na awarie czy warstwy API\@.
W~praktyce oznacza to, że użytkownik musi samodzielnie zbudować pozostałe elementy niezbędne do utrzymania jakości danych
oraz do ich wygodnego wykorzystania w~aplikacjach.
Projekt opisany w~pracy uzupełnia te braki poprzez wprowadzenie logicznego podziału na warstwę bieżącą i~archiwalną,
procesy ETL, interfejs GraphQL dla zapytań oraz kanały powiadomień.

\subsection{Wnioski z analizy porównawczej}
\label{subsec:wnioski-z-analizy-porownawczej}

Przedstawione projekty potwierdzają, że problem pozyskiwania i~wykorzystywania danych z~giełdy Binance jest szeroko
podejmowany, jednak najczęściej realizowany jest w~formie rozwiązań fragmentarycznych.
Jedne narzędzia skupiają się na wizualizacji, inne na wydajnym pobieraniu danych historycznych, a~jeszcze inne na samym
przetwarzaniu strumieniowym.
W~każdym z~przypadków widoczne są ograniczenia z~perspektywy docelowego zastosowania w~systemie analitycznym, tj.\ brak
spójnego potoku danych, brak mechanizmów retencji i~migracji między warstwami pamięci, brak wygodnego interfejsu zapytań
oraz brak funkcji powiadomień.

System zaproponowany w~niniejszej pracy stanowi próbę integracji tych rozproszonych funkcjonalności w~jedno rozwiązanie,
które umożliwia przetwarzanie danych w~czasie rzeczywistym, ich kontrolowaną archiwizację oraz późniejsze wykorzystanie
analityczne.
